\documentclass[12pt]{article}

\usepackage[utf8]{inputenc}
\usepackage{amsmath, amssymb}
\usepackage{geometry}
\usepackage{tikz}
\usetikzlibrary{arrows.meta}

\geometry{margin=1in}
\setlength{\parindent}{0pt}

\title{Effect of a Disability}
\author{ECON 300 -- Labour Economics}
\date{Winter 2026}

\begin{document}
\maketitle
\hrule



\vspace{1em}
\section*{General Framework (Both Panels)}

\begin{itemize}
\item The vertical axis measures \textbf{income} ($Y$).
\item The horizontal axis measures \textbf{leisure}.
\item The downward-sloping straight line represents the \textbf{budget constraint}.
\item Indifference curves represent preferences over income and leisure.
\item Optimal choices occur at points where an indifference curve is tangent to the budget constraint.
\item Hours of work are defined as total available time minus leisure.
\end{itemize}

\vspace{0.5em}

\section*{Panel (a): Reduced Work Time --- Partial and Full Disability}

\textbf{Key idea:} Disability reduces the individual’s \emph{maximum feasible working time}.

\begin{itemize}
\item Before disability:
  \begin{itemize}
  \item The individual can supply up to $H_f$ hours of work.
  \item The optimal choice is at point $E_o$, associated with indifference curve $U_o$.
  \item Hours worked are $H_o$.
  \end{itemize}

\item After disability:
  \begin{itemize}
  \item The individual’s maximum feasible work time is reduced.
  \item Under partial disability, maximum hours fall to $H_p$.
  \item Under full disability, feasible work time may fall further.
  \item The budget constraint becomes \textbf{shorter}, ending at $H_p$ instead of $H_f$.
  \end{itemize}

\item Because the old optimum $E_o$ is no longer feasible:
  \begin{itemize}
  \item The individual chooses a new optimum $E_p$.
  \item The new indifference curve $U_p$ lies below $U_o$, indicating lower utility.
  \item Hours worked fall from $H_o$ to $H_p$.
  \end{itemize}
\end{itemize}

\textbf{Economic interpretation:}
\begin{itemize}
\item Disability acts as a \textbf{constraint on the time endowment}.
\item Labour supply falls mechanically because the individual cannot physically work as many hours.
\item Preferences and wages are unchanged, but feasible choices are restricted.
\end{itemize}

\vspace{0.5em}

\section*{Panel (b): Increased Disutility of Labour}

\textbf{Key idea:} Disability increases the \emph{disutility of working} without reducing available time.

\begin{itemize}
\item The budget constraint remains unchanged:
  \begin{itemize}
  \item Wage rate is the same.
  \item Maximum available time is unchanged.
  \end{itemize}

\item Disability increases the physical or mental cost of working:
  \begin{itemize}
  \item Working becomes more painful or exhausting.
  \item Leisure becomes relatively more valuable.
  \end{itemize}

\item Preferences change:
  \begin{itemize}
  \item Indifference curves become \textbf{steeper}.
  \item Original indifference curve: $U_o$.
  \item New indifference curves after disability: $U_d'$ and $U_d''$.
  \end{itemize}

\item Optimal choice shifts:
  \begin{itemize}
  \item Tangency moves to the right.
  \item Leisure increases.
  \item Hours worked fall from $H_o$ to $H_d$.
  \end{itemize}
\end{itemize}

\textbf{Economic interpretation:}
\begin{itemize}
\item Disability affects \textbf{preferences}, not constraints.
\item Labour supply falls because the marginal disutility of work rises.
\item This is a behavioural response rather than a mechanical restriction.
\end{itemize}

\vspace{0.5em}

\section*{Comparison of the Two Mechanisms}

\begin{center}
\begin{tabular}{p{0.32\linewidth} p{0.3\linewidth} p{0.3\linewidth}}
\hline
 & \textbf{Panel (a)} & \textbf{Panel (b)} \\
\hline
What changes & Time constraint & Preferences \\
Budget constraint & Shorter & Unchanged \\
Indifference curves & Same & Steeper \\
Reason hours fall & Physical limit & Higher disutility \\
Nature of response & Mechanical + choice & Choice only \\
\hline
\end{tabular}
\end{center}

\vspace{0.5em}

\section*{Key Takeaway (Chapter 3)}

\begin{itemize}
\item Disability can reduce labour supply through two distinct channels:
  \begin{enumerate}
  \item Reduced capacity to work (time constraint).
  \item Increased disutility of working (preference change).
  \end{enumerate}
\item Observed reductions in work do not necessarily imply a reduced willingness to work.
\item Effective disability and income support policies must account for both mechanisms.
\end{itemize}

\end{document}
